\documentclass[11pt,twoside]{article}

\usepackage{fullpage}

\usepackage{epsf}
\usepackage{fancyhdr}
\usepackage{graphics}
\usepackage{graphicx}
\usepackage{psfrag}
\usepackage{microtype}
\usepackage{tikz}
\usepackage{subfigure}
\usepackage{algorithm}
\usepackage{algpseudocode}

\usepackage{color}

\usepackage{amsthm}
\usepackage{amsfonts}
\usepackage{amsmath}
\usepackage{amssymb,bbm}

\usepackage{natbib}

\usepackage{url}% for url's in bib
% for theorem hyperlink colors
\usepackage[colorlinks,linkcolor=magenta,citecolor=blue,pagebackref=true]{hyperref}
% \renewcommand*{\backref}[1]{\ifx#1\relax \else Page #1 \fi}
\renewcommand*{\backrefalt}[4]{%
    \ifcase #1 \footnotesize{(Not cited.)}%
    \or        \footnotesize{(Cited on page~#2.)}%
    \else      \footnotesize{(Cited on pages~#2.)}%
    \fi}

\long\def\comment#1{}
% for nice fractions
\usepackage{nicefrac}


% for adjust width
\usepackage{chngpage}

 \usepackage{tabularx}%

% to label enumerate
\usepackage{enumitem}
% Top and bottom rules for tables
\usepackage{booktabs}
% for captions
\usepackage{caption}


% for mathmakebox
\usepackage{mathtools}

%
\usepackage{fullpage}



%%%%%%%%%%%%%%%%%%%%%%%%%%%%%%%%%%%%%%%%%%%%%%%%%%%%%%%%%%%%%%%%%%%%%%%%%%

%% COMMENTING commands

\long\def\comment#1{}
\definecolor{battleshipgrey}{rgb}{0.52, 0.52, 0.51}
\definecolor{darkgray}{rgb}{0.66, 0.66, 0.66}
\definecolor{darkgreen}{rgb}{0.0, 0.2, 0.13}
\definecolor{darkspringgreen}{rgb}{0.09, 0.45, 0.27}
\definecolor{dukeblue}{rgb}{0.0, 0.0, 0.61}
\definecolor{olivedrab7}{rgb}{0.24, 0.2, 0.12}
\definecolor{darkblue}{rgb}{0.0, 0.0, 0.55}
\definecolor{darkscarlet}{rgb}{0.34, 0.01, 0.1}
\definecolor{candyapplered}{rgb}{1.0, 0.03, 0.0}
\definecolor{ao(english)}{rgb}{0.0, 0.5, 0.0}
\definecolor{applegreen}{rgb}{0.55, 0.71, 0.0}

\newcommand{\red}[1]{\textcolor{red}{#1}}

\newcommand{\datcomment}[1]{{\bf{{\blue{{Dat --- #1}}}}}}
\newcommand{\blue}[1]{\textcolor{blue}{#1}}
\newcommand{\green}[1]{\textcolor{green}{#1}}
\newcommand{\highlight}[1]{\textcolor{applegreen}{#1}}


\setlength{\textwidth}{\paperwidth}
\addtolength{\textwidth}{-6cm}
\setlength{\textheight}{\paperheight}
\addtolength{\textheight}{-4cm}
\addtolength{\textheight}{-1.1\headheight}
\addtolength{\textheight}{-\headsep}
\addtolength{\textheight}{-\footskip}
\setlength{\oddsidemargin}{0.5cm}
\setlength{\evensidemargin}{0.5cm}

\def\proof{\par\noindent{\em Proof.}}
\def\endproof{\hfill $\Box$ \vskip 0.4cm}

\renewcommand{\theHalgorithm}{\arabic{algorithm}}

\RequirePackage{natbib}
%\RequirePackage[OT1]{fontenc}
%\RequirePackage{amsthm,amsmath}
%\RequirePackage[numbers]{natbib}
%\RequirePackage[colorlinks,citecolor=blue,urlcolor=blue]{hyperref}
%\RequirePackage[dvips]{hyperref}
%\usepackage{natbib}

\newcommand{\myline}{1.5}
\newcommand{\trace}{\ensuremath{tr}}

\long\def\comment#1{}
\renewcommand\vec[1]{\ensuremath\boldsymbol{#1}}
\newcommand{\Sdd}{\ensuremath{S_{d_2}^{++}}}
\renewcommand{\trace}{\ensuremath{\operatorname{Tr}}}
\newcommand{\Ecal}{\ensuremath{\mathcal{E}}}
\newcommand{\Tcal}{\ensuremath{\mathcal{T}}}
\newcommand{\Ocal}{\ensuremath{\mathcal{O}}}
\newcommand{\Gcal}{\ensuremath{\mathcal{G}}}
\newcommand{\Xcal}{\ensuremath{\mathcal{X}}}
\newcommand{\Bcal}{\ensuremath{\mathcal{B}}}
\newcommand{\Acal}{\ensuremath{\mathcal{A}}}
\newcommand{\Ccal}{\ensuremath{\mathcal{C}}}

\newcommand{\Ycal}{\ensuremath{\mathcal{Y}}}
\newcommand{\Pcal}{\ensuremath{\mathcal{P}}}
\newcommand{\Ncal}{\ensuremath{\mathcal{N}}}
\newcommand{\Kcal}{\ensuremath{\mathcal{K}}}
\newcommand{\Rbb}{\ensuremath{\mathbb{R}}}
\newcommand{\Ebb}{\ensuremath{\mathbb{E}}}
\newcommand{\Nbb}{\ensuremath{\mathbb{N}}}
\newcommand{\Hcal}{\ensuremath{\mathcal{H}}}
\newcommand{\Mcal}{\ensuremath{\mathcal{M}}}
% \newcommand{\emptyset}{\ensuremath{\varphi}}

\newcommand{\Hbb}{\ensuremath{\mathbb{H}}}
\newcommand{\Qbb}{\ensuremath{\mathbb{Q}}}
\newcommand{\Bbbb}{\ensuremath{\mathbb{B}}}
\newcommand{\supp}{\text{supp}}
\newcommand{\knowfunmix}{\ensuremath{h_{2}}}
\newcommand{\knowfunsec}{\ensuremath{h_{1}}}
\newcommand{\underfun}{\ensuremath{g_{*}}}
\newcommand{\mleprob}{\ensuremath{\widehat{G}_{n}}}
\newcommand{\mleweight}{\ensuremath{\widehat{\lambda^{0}}_{n}}}
\newcommand{\sureprob}{\ensuremath{\overline{G}_{n}}}
\newcommand{\sureweight}{\ensuremath{\overline{\lambda^{0}}_{n}}}
\newcommand{\ambispace}{\ensuremath{\mathcal{G}}}
\newcommand{\overlapset}{\ensuremath{\mathcal{A}}}
\newcommand{\mydefn}{\ensuremath{: =}}
\newcommand{\gaussden}{\ensuremath{\phi}}
\newcommand{\extraset}{\ensuremath{\mathcal{B}}}
\newcommand{\extraind}{\ensuremath{\mathcal{I}}}
\newcommand{\sumind}{\ensuremath{\mathcal{S}}}
\newcommand{\bigO}{\ensuremath{\mathcal{O}}}
\newcommand{\smallO}{\ensuremath{o}}
\newcommand{\diff}{\ensuremath{\mathcal{T}}}
\newcommand{\distrong}{\ensuremath{\mathcal{D}_{\text{strong}}}}
\newcommand{\disgauss}{\ensuremath{\mathcal{D}_{\text{gauss}}}}
\DeclareMathOperator*{\argmax}{arg\,max}
\DeclareMathOperator*{\argmin}{arg\,min}

\newcommand{\Xbf}{\ensuremath{\mathbf{X}}}
\newcommand{\Rbf}{\ensuremath{\mathbf{R}}}
\newcommand{\Qbf}{\ensuremath{\mathbf{Q}}}
\newcommand{\Pbf}{\ensuremath{\mathbf{P}}}
\newcommand{\Gbf}{\ensuremath{\mathbf{G}}}
\newcommand{\fbf}{\ensuremath{\mathbf{f}}}
\newcommand{\ubf}{\ensuremath{\mathbf{u}}}
\newcommand{\Zbf}{\ensuremath{\mathbf{Z}}}


%\newcommand{\disgaussmod}{\ensuremath{\bar{\mathcal{D}}_{\text{gauss}}}}
\newcommand{\diskew}{\ensuremath{\mathcal{D}_{\text{skew}}}}

%%% Brackets
\newcommand{\brackets}[1]{\left[ #1 \right]}
\newcommand{\parenth}[1]{\left( #1 \right)}
\newcommand{\bigparenth}[1]{\big( #1 \big)}
\newcommand{\biggparenth}[1]{\bigg( #1 \bigg)}
\newcommand{\braces}[1]{\left\{ #1 \right \}}
\newcommand{\abss}[1]{\left| #1 \right |}
\newcommand{\angles}[1]{\left\langle #1 \right \rangle}
\newcommand{\tp}{^\top}

% Some vector/matrix norms
\newcommand{\matnorm}[3]{|\!|\!| #1 | \! | \!|_{{#2}, {#3}}}
\newcommand{\matsnorm}[2]{|\!|\!| #1 | \! | \!|_{{#2}}}
\newcommand{\vecnorm}[2]{\left\| #1\right\|_{#2}}
\newcommand{\enorm}[1]{\vecnorm{#1}{2}} % euclidean norm
\newcommand{\nucnorm}[1]{\ensuremath{\matsnorm{#1}{\footnotesize{\mbox{nuc}}}}}
\newcommand{\fronorm}[1]{\ensuremath{\matsnorm{#1}{\footnotesize{\mbox{fro}}}}}
\newcommand{\opnorm}[1]{\ensuremath{\matsnorm{#1}{\tiny{\mbox{op}}}}}
\newcommand{\norm}[1]{\left\|#1\right\|}
\newcommand{\innprod}[1]{\left\langle#1\right\rangle}
\comment{
% text margin
\setlength{\topmargin}{0 in}
\setlength{\textwidth}{5.5 in}
\setlength{\textheight}{8 in}
\setlength{\oddsidemargin}{0.5 in}
}

\theoremstyle{plain}

\newtheorem{theorem}{Theorem}

\numberwithin{theorem}{section}

\newtheorem{proposition}{Proposition}

\numberwithin{proposition}{section}

\newtheorem{lemma}{Lemma}

\numberwithin{lemma}{section}

\newtheorem{definition}{Definition}

\numberwithin{definition}{section}

\newtheorem{condition}{Condition}

\numberwithin{condition}{section}

\newtheorem{problem}{Problem}

\numberwithin{problem}{section}

\newtheorem{corollary}{Corollary}

\numberwithin{corollary}{section}

\newtheorem{assumption}{Assumption}

\numberwithin{assumption}{section}

\newtheorem{example}{Example}

\numberwithin{example}{section}

\newtheorem{conjecture}{Conjecture}

\numberwithin{conjecture}{section}

%\newtheorem{example}{Example}


\theoremstyle{definition}

\newtheorem{remark}{Remark}

\numberwithin{remark}{section}

\newcommand{\floor}[1]{\lfloor #1 \rfloor}

\begin{document}

\begin{center}

{\bf{\LARGE{Mixture of Sub-manifolds Non-negative Matrix Factorization (MoS-NMF)}}}

\vspace*{.2in}
% {\large{
% \begin{tabular}{cc}
% Dat Do$^\diamond$, Sunrit Chakraborty$^\star$, Linh Do$^\ast$,\\ Jonathan Terhorst$^\dagger$, Matthew Stephens$^{\diamond, \ddagger}$
% \end{tabular}
% }}

% \vspace*{.2in}

% \begin{tabular}{c}
% $^\diamond$ Department of Statistics, University of Chicago, Chicago, IL \\
% $^\star$Department of Statistics, Duke University, Durham, NC \\ 
% $^\ast$Duke\\ 
% $^\dagger$ Department of Statistics, University of Michigan, Ann Arbor MI \\ 
% $^\ddagger$Department of Human Genetics, University of Chicago, Chicago, IL \\ 
% \end{tabular}

\vspace*{.2in}
\today
\vspace*{.2in}
\end{center}

\begin{abstract}
    
\end{abstract}

\section{Introduction}
Considering Poisson-NMF likelihood:
$$Y_{nm} \sim \mathrm{Poi}\left(\sum_{k=1}^{K} L_{nk} F_{km}\right),$$
for count data $(Y_{nm})_{n, m= 1}^{N, M}$, where $Y_{n m}$ counts how many times item $m$ appears in data/cell/document $n$. $L \in \mathbb{R}_{+}^{N\times K}$ capture the loadings/weights and $F \in 
\mathbb{R}_{+}^{K\times M}$ are $K$ factors/parts. \textcolor{blue}{$L$ is the observation's specific score.}

\paragraph{Modeling goal.} We want to model that each $Y_n$ is a combination of not all $K$ parts, but only a few of them. For example, even if we fit $K = 50$ topics to the copura, it makes sense to assume that each document can only discuss a maximum of 3-5 topics. Similarly, in music, each song can only be mixed of at most 2-3 genres. So we want to model each loading $L_n \in \Rbb_{+}^{K}$ to have at most $D$ non-zero effects. There are many ways to choose a combination of $D$ topics from $K$ topics, and we call each such way a "motif". Hence, we have $S$ sub-manifolds, each of which has a dimension of at most $D$ in $\Rbb_+^K$, which is controlled by parameter $\square_s := (\gamma_s, \alpha^{0}_s, \lambda^{0}_s)$ for $\gamma_s = (\gamma_{sd})_{d = 1}^{D} \in [K]^{D}$ indicates which $D$ topics are chosen and $(\alpha^{0}_{sd}, \lambda^{0}_{sd})_{s, d = 1}^{S, D}$ controls the randomness of the effect size. In other word, given the sub-manifold $Z_n$ of the loading $L_n$, we have
\begin{equation}
    (L_n | Z_n = s) \sim \sum_{d = 1}^{D} e_{\gamma_{sd}} \beta_{nd}, \quad \beta_{nd} \sim \mathrm{Gam}(\alpha^{0}_{sd}, \lambda^{0}_{sd}),
\end{equation}
where $e_{k}$ is the vector of 0 except a single 1 at $k$-th element. We call each of these distributions $\mathrm{Gam}(\square_s)$. Thus each $L_n$ is modeled as a mixture of them:
\begin{equation}
    L_n \overset{iid}{\sim} \sum_{s=1}^{S} \pi_s \mathrm{Gam}(\square_s).
\end{equation}
We finish the model by putting uniform prior on each $\gamma_s$:
\begin{equation}
    \gamma_{sd} \overset{iid}{\sim} \mathrm{Unif}([K]), \quad s\in [S], d \in [D].
\end{equation}
The observation is 
\begin{equation}
    Y_{nm} | L_n, F \overset{ind.}{\sim} \mathrm{Poi}(\sum_{k=1}^{K} L_{nk} F_{km}), \quad n\in [N], m\in [M].
\end{equation}
We will make inference of $(\pi_{s})_{s=1}^{S}, (\alpha^{0}_{sd}, \lambda^{0}_{sd})_{s, d =1}^{S, D}$ and $F$ by mean of Maximum Likelihood, and make inference of indicator $\gamma$ and effect size $(\beta_{nd})_{N, D}$ by its posterior (given the MLE of other parameters). So this is an Empirical Bayes procedure. The reason why we want to use the Bayesian formula for the indicator $\gamma$ is that we think it can quantify the uncertainty if there are some very similar $F_k$'s (detect over-fitting). Let's see if it works as intended.

\paragraph{The whole model.} Hence, the whole model is as follows:
\begin{equation}\label{eq:whole-model}
\begin{cases}
\gamma_{sd} \overset{iid}{\sim} \mathrm{Unif}([K]), & s\in [S], d\in [D]\\
Z_{n} | \pi \overset{iid}{\sim} \mathrm{Cat}(\pi) \in [S], & n\in [N]\\
\beta_{nd} | (\alpha^{0}, \lambda^{0}, Z_n = s) \overset{ind.}{\sim} \mathrm{Gam}(\alpha^{0}_{sd}, \lambda^{0}_{sd}), & n\in[N], s\in [S], d\in [D]\\
Y_{nm} | (\gamma, \beta, Z_n = s) \overset{ind.}{\sim} \mathrm{Poi}(\sum_{d = 1}^{D} \beta_{nd} F_{\gamma_{sd}m}), & n\in [N], m\in [M],
\end{cases}
\end{equation}
with latent variables $(Z_n)_{n=1}^{N}, (\gamma_{sd})_{s, d = 1}^{S, D}, (\beta_{nd})_{n, d = 1}^{N, D}$ and parameters $(\pi_s)_{s=1}^{S}, (\alpha^{0}_{sd})_{s, d = 1}^{S, D}, (\lambda^{0}_{sd})_{s, d=1}^{S, D}$, and $F\in \Rbb_{+}^{K\times M}$. For the identifiability reason, we will assume that $\sum_m F_{km} = 1$ for all $k\in [K]$.

\section{Preliminaries}
\subsection{Poisson SuSiE}
MoS-NMF is a complex hierarchical model that consists of (1) Poisson-SuSiE components and (2) mixture modeling. In this section, we present a backbone SuSiE-Poisson regression model for one observation and its inference algorithm, which helps build intuition and can later be incorporated into the inference of the MoS-NMF model. Given an observation $y = (y_1, \dots, y_m) \in \Nbb^{M}$ and factors $F \in \mathbb{R}_{+}^{K\times M}$, consider the sparse regression model 
$$y \sim \mathrm{Poi}(l F),$$
where $l \in \Rbb_{+}^{K}$, and the Poisson distribution is independent across $M$ dimensions, i.e.,
$$y_m \overset{ind.}{\sim} \mathrm{Poi}(\sum_{k = 1}^{K} l_k F_{km}), \quad m\in [M]. $$
We want to model the sparsity of $l\in \Rbb^{K}$ using the Sum of Single Effect model with at most $D$ non-zero effect:
$$l = \sum_{d = 1}^{D} l^{(d)} = \sum_{d = 1}^{D} e_{\gamma_{d}} \beta_{d},$$
where $\gamma_{d} \sim \mathrm{Unif}([K])$ (sparsity index indicator) and $\beta_{d} \sim \mathrm{Gamma}(\alpha^{0}_{d}, \lambda^{0}_{d})$ a priori (Gamma distribution is parameterized using shape-rate parameter here). We want to make inference of the posterior $(\gamma_{d}, \beta_{d})_{d = 1}^{D} | y$. 
\paragraph{Poisson Single Effect Regression.} When $D = 1$, we can explicitly calculate the posterior
\begin{align*}
p(y | \gamma_{1} = k) & = \int p(x | \gamma_{1} = k, \beta_1) \mathrm{Gamma}(\beta_{1} | \alpha^{0}_1, \lambda^{0}_1) d\beta_{1} \\
& = \int \prod_{m=1}^{M} \mathrm{Poi}(y_m | \beta_1 F_{km}) \mathrm{Gamma}(\beta_{1} | \alpha^{0}_1, \lambda^{0}_1) d\beta_{1} \\
& = \int_{\beta_{1}} \dfrac{e^{-\beta_{1} \sum_m F_{km}}\prod_m (\beta_{1} F_{km})^{y_{m}}}{\prod_{m} y_{m} !} \dfrac{(\lambda^{0}_1)^{\alpha^{0}_1}}{\Gamma(\alpha^{0}_1)} e^{-\beta_{1} \lambda^{0}_1} \beta_1^{\alpha^{0}_{1} - 1} d \beta_1\\
& = \dfrac{\prod_{m=1}^{M} F_{km}^{y_{m}}}{\prod_{m=1}^{M} y_{m}!} \dfrac{(\lambda^{0}_1)^{\alpha^{0}_1}}{\Gamma(\alpha^{0}_1)} \dfrac{\Gamma(\sum_m y_{m} + \alpha^{0}_1)}{(\sum_{m} F_{km} + \lambda^{0}_1)^{\sum_m y_{m} + \alpha^{0}_1}}. 
\end{align*}
Hence,
\begin{align}
    p(\gamma_1 = k | y) \propto \dfrac{\prod_{m=1}^{M} F_{km}^{y_{m}}}{(\sum_{m} F_{km} + \lambda^{0}_1)^{\sum_m y_{m} + \alpha^{0}_1}}, \quad k \in [K].
\end{align}
Moreover, by the Poisson-Gamma conjugacy, it is straightforward that:
\begin{align}
    p(\beta_1 | y, \gamma_1 = k) \sim \mathrm{Gam}\left(\sum_m y_m + \alpha^{0}_1, \sum_{m} F_{km} + \lambda^{0}_1\right), \quad k \in [K].
\end{align}

Hence, the posterior distribution $p(\gamma_1, \beta_1 | x)$ can be summarized by $3K$ non-negative number. (This point can be simplified by noticing the posterior of $\beta_1$.) For the sake of future argument, we also note that this exact posterior distribution has the variational characterization that
\begin{equation}
    p(\gamma_1, \beta_1 | y) = \argmax_{q(\gamma_1, \beta_1)} \Ebb_q \left[\sum_m \log \mathrm{Poi}(y_m | \beta_1 F_{\gamma_1 m}) \right] - \mathrm{KL}(q(\gamma_1, \beta_1) \| p(\gamma_1, \beta_1)),
\end{equation}
where $\mathrm{Poi}(y_m | \beta_1 F_{\gamma_1 m}) = y_m \log(\beta_1) + \log(F_{\gamma_1 m}) - F_{\gamma_1 m } \beta_1 - \log (y_m!)$


\paragraph{Poisson SuSiE.} The posterior of $(\gamma_{d}, \beta_{d})_{d = 1}^{D}$ does not have a simple closed-form expression, so we use the same Variational approximation as in Gaussian SuSiE~\cite{Wang2020-cr}:
\begin{equation}
    p(\gamma, \beta | y) \approx \prod_{d = 1}^{D} q_d(\gamma_d, \beta_d) = \prod_{d = 1}^{D} (q(\gamma_d) q(\beta_d | \gamma_d)),
\end{equation}
and find the optimal variational distribution by maximizing the ELBO:
\begin{align*}
\Ebb_{\prod_d q_d} \left[\sum_{m} \log \mathrm{Poi} \bigg(y_m \mid \sum_{d=1}^{D} \beta_d F_{\gamma_d m}\bigg)\right] - \sum_{d=1}^{D} \mathrm{KL}(q_d(\gamma_d, \beta_d) \| p(\gamma_d, \beta_d)). 
\end{align*}
Note that in the log-likelihood: 
\begin{align*}
    \log \mathrm{Poi} \bigg(y_m \mid \sum_{d=1}^{D} \beta_d F_{\gamma_d m}\bigg) = y_m \log\bigg(\sum_{d=1}^{D} \beta_d F_{\gamma_d m}\bigg) -  \sum_{d=1}^{D} \beta_d F_{\gamma_d m} - \log (y_m!),
\end{align*}
the third term is a constant, the expectation of the second term can be calculated simply as
\begin{equation}
\Ebb_q \sum_{d=1}^{D} \beta_d F_{\gamma_d m} = \sum_{d=1}^{D} \sum_{k=1}^{K} p(\gamma_d = k)  \Ebb [\beta_d | \gamma_d = k] F_{km},
\end{equation}
while the first term does not have a simple closed-form expectation. So that we utilize a well-known trick in latent variable model to introduce new auxiliary variables $(\xi_{md})_{m, d=1}^{M, D}$ and use the lower bound:
\begin{equation}
\log\bigg(\sum_{d=1}^{D} \beta_d F_{\gamma_d m}\bigg) \geq \sum_{d=1}^{D} \xi_{md} \left(\log \left(\dfrac{\beta_d F_{\gamma_d m}}{\xi_{md}}\right) \right),
\end{equation}
where $\sum_{d} \xi_{md} = 1$ for all $m\in [M]$. 
So that we are maximizing the following Evidence Lower Bound:
\begin{align}
\mathcal{E}(q) & = \sum_{d=1}^{D} \Ebb_{q_d(\gamma_d, \beta_d)} \left[\sum_m \left(\xi_{md} y_m (\log(\beta_d) + \log(F_{\gamma_d m})) - \beta_d F_{\gamma_d m} \right)\right] \nonumber \\
& \,\, - \sum_{d=1}^{D}  \mathrm{KL}(q(\gamma_d, \beta_d)\| p(\gamma_d, \beta_d)) - \sum_{m, d} y_m \xi_{md} \log(\xi_{md}).  
\end{align}

\begin{algorithm}
\caption{SuSiE-Poisson algorithm}\label{alg:susie-poisson}
\begin{algorithmic}
\For{Outer loop}
    \For {$d$ from 1 to $D$}
    \State Update auxiliary variable $\xi_{md} \gets  \text{Norm-Exp}(\Ebb_{q(\gamma_{d}, \beta_{d})} (\log \beta_{d} + \log(F_{\gamma_{d}m}))) \quad \forall m\in [M], d\in [D].$ \Comment{(Gauss-Seidel style update)}
    \State Update $\alpha_{dk} \gets \sum_m \xi_{md} y_m + \alpha^{0}_d$ \Comment{(Variational parameters)}
    \State Update $\lambda_{dk} \gets \sum_m F_{km} + \lambda^{0}_d$ \Comment{(Variational parameters)}
    \State Assign $q(\beta_d | \gamma_d = k) \equiv \mathrm{Gam}(\alpha_{dk}, \lambda_{dk})$ \Comment{(Variational distributions)}
    \State Calculate $\mathrm{logprob}_{dk} = \sum_m \xi_{md} y_m \log(F_{km}) - (\sum_m \xi_{md} y_m + \alpha^{0}_d)\log(\sum_{m} F_{km} + \lambda^{0}_d)$
    \State Update $q(\gamma_d) = (\bar{\gamma}_{dk})_{k=1}^{K} \gets \text{Norm-Exp}(\mathrm{logprob}_{d1}, \dots, \mathrm{logprob}_{dK}) $ \Comment{(PIP)}
    \State Update $(\alpha_d^0, \lambda_d^0)$ \Comment{(Optional; VEB)}
    
    \EndFor
    
\EndFor
\end{algorithmic}
\end{algorithm}

\paragraph{ELBO.}
Let $q(\gamma_d = k) \equiv \bar{\gamma}_{dk}$ and $q(\beta_d | \gamma_d = k) \equiv \mathrm{Gam}(\alpha_{dk}, \lambda_{dk})$, where 
$$\alpha_{dk} = \sum_{m} \xi_{md} y_{m} + \alpha_d^0, \quad \lambda_{dk} = \sum_{m} F_{km} + \lambda_d^0.$$
We have
\begin{equation}
\Ebb_{q(\beta_d | \gamma_d = k)} [\log(\beta_{d})] = \psi(\alpha_{dk}) - \log(\lambda_{dk}),  \quad \Ebb_{q(\beta_d, \gamma_d)} [\log(\beta_{d})] = \sum_k \bar{\gamma}_{dk}[\psi(\alpha_{dk}) - \log(\lambda_{dk})],  
\end{equation}
and
\begin{align}
\mathrm{KL}(q_d \| p_d) & = \sum_{k=1}^{K} \bar{\gamma}_{dk}(\log(\bar{\gamma}_{dk}) - \log(\tfrac{1}{K})) \nonumber \\
& + \sum_{k=1}^{K} \bar{\gamma}_{dk} \left[(\alpha_{dk} - \alpha_d^0) \psi(\alpha_{dk}) + (\log(\lambda_{dk}) - \log(\lambda_d^0))\alpha^0_d - \log \dfrac{\Gamma(\alpha_{dk})}{\Gamma(\alpha_{d}^0)} - (\lambda_{dk} - \lambda_d^0) \dfrac{\alpha_{dk}}{\lambda_{dk}} \right].   
\end{align}

\begin{remark}
Our experiments show that updating $\xi$ every $d$ (Gauss-Seidel style) helps the algorithm escape local optima much more frequently than updating $\xi$ only once per outer loop (Jacobi style).
\end{remark}

\subsection{Mixture modeling}

\section{Model fitting in MoS-NMF}
We approximate the posterior as: 
$$p(Z, \gamma, \beta | X) \approx q(Z, \gamma, \beta) =  \prod_{n=1}^{N} q(Z_n) \prod_{s, d = 1}^{S, D} q(\gamma_{sd}) \prod_{n, d} q(\beta_{nd} | Z_n, \gamma_{Z_n d}).$$
\textcolor{red}{Should I factor $\beta$, or the variational lower bound implies it?}
Note that both $Z$ and $\gamma$ are discrete. $Z_n$ takes value in $[S]$, and $\gamma$ in $[K]$. Write 
$$\tilde{z}_{ns} := q(Z_n = s), \quad \tilde{\gamma}_{sdk} := q(\gamma_{sd} = k). $$
Also notice that the likelihood can be decomposed as:
$$\log p(X | Z, \gamma, \beta) = \sum_{n=1}^{N} \log p(Y_n | Z_n, \gamma, \beta_n), \quad \log p(\beta | Z, \alpha^{0}, \lambda^{0}) = \sum_{n=1}^{N} p(\beta_n | Z_n, \alpha^{0}, \lambda^{0}),$$
where
$$\log p(Y_n | Z_n = s, \gamma, \beta_n) = \sum_{m} \log \mathrm{Poi}(Y_{nm} | \sum_{d=1}^{D} \beta_{nd} F_{\gamma_{sd}m}).$$
and 
$$\log p(\beta_n | Z_n = s, \alpha^{0}, \lambda^{0}) = \sum_{d = 1}^{D} \log\mathrm{Gam}(\beta_{nd} | \alpha^{0}_{sd}, \lambda^{0}_{sd}).$$
Using variational lower bound:
\begin{align*}
p(X) & \geq \Ebb_{q(Z, \gamma, \beta)} \big\{\log p(X|  Z, \gamma, \beta) + \log p(\beta | \alpha^{0}, \lambda^{0}, Z) + \log p(Z | \pi) + \log p(\gamma) \\
& \hspace{2cm} - \log q(Z) - \log q(\gamma) - \log q(\beta | Z, \gamma) \big\}\\
& = \sum_{n=1}^{N} \sum_{s=1}^{S} \tilde{z}_{ns} \Ebb_{q(\gamma) q(\beta | Z, \gamma)} \big\{\log p(Y_n| Z_n = s, \gamma, \beta) + \log p(\beta_n |Z_n = s, \alpha^{0}, \lambda^{0}) \\
& \hspace{1cm}  + \log p(Z_n = s| \pi) + \log p(\gamma) - \log q(Z_n = s) - \log q(\gamma) - \log q(\beta_n | Z_n = s, \gamma) \big\}\\
& = \sum_{n=1}^{N} \sum_{s=1}^{S} \tilde{z}_{ns} \Ebb_{q(\gamma) q(\beta | Z, \gamma)} \bigg\{ \sum_{m=1}^{M} \log \mathrm{Poi}\big(Y_{nm}| \sum_{d} \beta_{nd} F_{\gamma_{sd}m} \big) + \sum_{d} \log \mathrm{Gam}(\beta_{nd} |\alpha^{0}_{sd}, \lambda^{0}_{sd}) \\
& \hspace{1cm}  + \log p(\gamma) - \log q(\gamma) - \log q(\beta_n | Z_n = s, \gamma) \bigg\} + \sum_{n=1}^{N} \sum_{s=1}^{S} \tilde{z}_{ns} \big(\log(\pi_s) - \tilde{z}_{ns} \big)
\end{align*}
To efficiently estimate the expectation of the log-likelihood with respect to $q(\gamma)$, we further use a lower bound with auxiliary variables $(\xi_{nmd})_{n, m, d}^{N, M, D}$:
\begin{align}
& \log \mathrm{Poi}\big(Y_{nm}| \sum_{d} \beta_{nd} F_{\gamma_{sd}m} \big)  = - \sum_{d} \beta_{nd} F_{\gamma_{sd}m} + Y_{nm}\log\left(\sum_{d} \beta_{nd} F_{\gamma_{sd}m} \right) - \log (Y_{nm}!) \nonumber\\
&\geq \sum_{d} \left(- \beta_{nd} F_{\gamma_{sd} m} + Y_{nm} \xi_{nmd} \log \left(\dfrac{\beta_{nd} F_{\gamma_{sd}m}}{\xi_{nmd}}\right) \right) - \log (Y_{nm}!)
\end{align}
Ignoring the constant terms $\log (Y_{nm}!)$ and the terms $\log p(\gamma)$ (uniform distributions), since now the log-likelihood's lower bound can be written as a sum over $d$, we can now write the lower bound as:
\begin{align*}
p(X) & \geq \sum_{n, s = 1}^{N, S} \tilde{z}_{ns} \Ebb_{q(\gamma) q(\beta | Z, \gamma)} \bigg\{ \sum_{d, m = 1}^{D, M} \left(- \beta_{nd} F_{\gamma_{sd} m} + Y_{nm} \xi_{nmd} \log \left(\dfrac{\beta_{nd} F_{\gamma_{sd}m}}{\xi_{nmd}}\right) \right)  \\
& \hspace{1cm} + \sum_{d} \log \mathrm{Gam}(\beta_{nd} |\alpha^{0}_{sd}, \lambda^{0}_{sd})  + \log p(\gamma) - \log q(\gamma) - \log q(\beta_n | Z_n = s, \gamma) \bigg\} \\
& \hspace{1cm} + \sum_{n=1}^{N} \sum_{s=1}^{S} \tilde{z}_{ns} \big(\log(\pi_s) - \tilde{z}_{ns} \big)\\
& \geq \sum_{n, s, d, k = 1}^{N, S, D, K} \tilde{z}_{ns} \tilde{\gamma}_{sdk} \Ebb_{q(\beta | Z, \gamma)} \bigg\{ \sum_{m = 1}^{M} \left(- \beta_{nd} F_{k m} + Y_{nm} \xi_{nmd} \log \left(\dfrac{\beta_{nd} F_{km}}{\xi_{nmd}}\right) \right) 
\end{align*}


\begin{lemma}\label{lem:log_evidence}
For summary statistics $(T^{1}_{nd}, T^{\mathrm{log}}_{nd}, T^{0}_{nd})$ we have
\begin{align*}
& \max_{q} \Ebb_{q(\beta_{nd})} \left[ - T^{1}_{nd} \beta_{nd}  + (T^{\mathrm{log}}_{nd} - 1) \log(\beta_{nd}) + T^0_{nd} - \log q(\beta_{nd}) \right] \\
& = T_{nd}^0 - \log(T_{nd}^1) T^{\mathrm{log}}_{nd} + \log \Gamma(T_{nd}^{\mathrm{log}}).
\end{align*}
\end{lemma}
\begin{proof}[Proof of Lemma~\ref{lem:log_evidence}]
By means of the variational principle, 
$$q(\beta_{nd}) \propto \beta_{nd}^{(T^{\mathrm{log}}_{nd} - 1)} e^{- \beta_{nd} T^1_{nd}}, $$
which implies $q(\beta_{nd}) = \mathrm{Gam}(T^{\mathrm{log}}_{nd}, T^{1}_{nd})$. Hence,
\begin{equation}
\log q(\beta_{nd}) = (T^{\mathrm{log}}_{nd} - 1) \log(\beta_{nd}) - T^1_{nd} \beta_{nd} + \log(T_{nd}^{1}) T_{nd}^{\mathrm{log}} - \log \Gamma(T_{nd}^{\mathrm{log}}). 
\end{equation}
Plug this in, and we have the desire equation.
\end{proof}

\section*{Thoughts}
To fix the identifiability issue of the scales, we can either choose to
\begin{enumerate}
    \item Fix the scale of $F$, i.e., force $\sum_m F_{km} = 1$ for all $k$ and learn $\lambda^{0}_{sd}$ for all $s, d$
    \item Let $F$ be free and fix $\lambda^{0}_{sd} = 1$ for all $s, d$. 
\end{enumerate}
In the case that the modeling assumption of Gamma distribution does not hold, I think the first strategy leads to over-specifying $S$, while the second strategy leads to over-specifying $K$. I like the first one better. 



\bibliography{sparsity}
\bibliographystyle{apalike}




\end{document}
\section{Mixture of Manifolds Poisson NMF}
\paragraph{Poisson Sum of Single Effect model (Poi-SuSiE).} 
We are using a few lower bounds. First, 
\begin{equation}
    p(X) \geq \sum_{n, s} \omega_{ns}\left(\log(\pi_s) + \log (p(Y_n | \square_s)) - \log(\omega_{ns}) \right),
\end{equation}
for $\sum_s \omega_{ns} = 1$ for all $n\in [N]$. Now gotta lower bound
$$\sum_{n, s} \omega_{ns} \log (p(Y_n | \square_s).$$

Of course, we are usingthe  E-M algorithm for the mixture model:
\begin{enumerate}
    \item E-step: Compute 
    \begin{equation}
    \omega_{ns} = \dfrac{\pi_s p(Y_n | \square_s)}{\sum_{s'}\pi_{s'} p(Y_n | \square_{s'})}      
    \end{equation}
    \item Optimize
    $$\sum_{n, s=1}^{N, S} \omega_{ns} \log p(Y_n | \square_s, F) $$
\end{enumerate}



